\section{Evaluation Methodology}

The evaluation separates belief accuracy, evidence allocation, agent value,
and final deployment validity rather than reporting only an energy predictor.

\subsection{Research Questions}

\textbf{RQ1 -- Cross-fidelity calibration:} Do L0--L3 observations yield
calibrated predictions of unseen L4 energy, latency, and SLO feasibility, and
how many sparse scale probes repair an unseen-node-count shift?

\textbf{RQ2 -- Decision efficiency:} At equal real GPU-hours, how quickly does
\ipig\ recover the measured SLO-feasible Pareto set relative to random,
single-fidelity, multi-fidelity, fixed-ladder, and simulator-only search?

\textbf{RQ3 -- Agent contribution:} Does semantic subgoal selection and failure
repair improve success, cost, or reliability beyond an optimizer-only fixed
state machine, and does deterministic acquisition improve on an LLM-only agent?

\textbf{RQ4 -- Verified benefit:} Do L4 runs confirm lower energy than runtime
defaults and the strongest optimizer baseline at matched TTFT/TPOT SLOs?

\subsection{Campaign and Ground Truth}

Table~\ref{tab:evaluation-matrix} defines a campaign sized for dense evidence at
cheap levels and sparse evidence at expensive levels.  Llama-3.1-8B supplies
single-GPU phase coverage; Llama-3.1-70B supplies one-node and two-node
parallelism.  The primary runtime is vLLM, with one cross-runtime or precision
slice when allocation permits.  Prefill-heavy and decode-heavy token buckets
are complemented by Poisson/bursty arrival traces for L4.

\begin{table*}[t]
\centering
\caption{Minimum evaluation campaign.  Dense sampling is restricted to cheap
levels; multi-node measurements are selected sparsely and verified at L4.}
\label{tab:evaluation-matrix}
\scriptsize
\setlength{\tabcolsep}{4pt}
\renewcommand{\arraystretch}{1.1}
\begin{tabular}{p{0.12\textwidth}p{0.33\textwidth}p{0.26\textwidth}p{0.20\textwidth}}
\toprule
Dimension & Core campaign & Cross-scale / stress slice & Purpose \\
\midrule
Models & Llama-3.1-8B and 70B & Optional Qwen2.5-32B transfer slice & Device- and node-scale behavior \\
Scale & 1 GPU and 2/4/8 GPUs on one node & Two nodes and target arrival trace & L1--L4 fidelity \\
Hardware & One primary NVIDIA H100/H200-class cluster & Second GPU/interconnect regime if available & Topology transfer \\
Engines & vLLM primary & One TensorRT-LLM or precision slice & Runtime sensitivity \\
Workloads & Prefill-heavy and decode-heavy token buckets & Poisson/bursty traces and long context & Phase and queue stress \\
Knobs & batch/token limits, chunked prefill, TP, power cap & placement, node count, BF16/FP8 where stable & Energy/SLO Pareto space \\
Held out & random configurations and one token/batch region & one GPU/node count and time-split trace & Interpolation and extrapolation \\
\bottomrule
\end{tabular}
\end{table*}


Ground truth is collected by \bench\ with synchronized phase timestamps and
GPU, CPU/DRAM, and node telemetry.  Broad L1--L3 points use at least three
post-warmup repetitions; final L4 comparisons use at least five matched-trace
repetitions and bootstrap confidence intervals.  Run order is randomized to
reduce thermal and temporal bias.  PDU or external-meter measurements are used
on a representative subset when available.

We construct replay environments from dense sweeps.  Hidden rows act as L4
truth, while measured subsets, simulator outputs, and explicit cross-scale
masks expose L0--L3 evidence with actual or normalized GPU-hour cost.  Splits
are fixed before fitting, so history retrieval cannot reveal hidden L4 rows.
The protocols hold out (i) random configurations, (ii) a token/batch region,
(iii) one GPU or node count, and (iv) a model/runtime slice.  Replay supports
many controlled policy episodes and an exhaustive oracle; top recommendations
are then submitted through the unchanged ClusterExecutor.

\subsection{Baselines, Ablations, and Metrics}

Search curves compare random search, single-fidelity qNEHVI
\cite{daulton2021qnehvi}, MF-OSEMO \cite{belakaria2020mfosemo}, constrained
MF-JESMOC \cite{fernandezsanchez2025mfjesmoc}, a fixed L0--L4 ladder, and a
simulator-only policy.  Every method receives identical initial evidence,
candidate pools, fidelity costs, random seeds, and L4 reserve.  Deployment
comparisons use the runtime default, the strongest optimizer recommendation,
and \system.  Agent ablations compare the full hybrid agent with an
optimizer-only fixed state machine and an LLM-only tool user.  Failure episodes
cover unsupported configurations, out-of-memory, preemption, telemetry loss,
and budget exhaustion.

Prediction metrics are MAPE/MAE, rank correlation, interval coverage, and SLO
calibration.  Decision metrics are feasible Pareto precision/recall, normalized
hypervolume regret, GPU-hours to a target regret, and L4 query count.
Deployment metrics include J/token, J/request, TTFT/TPOT percentiles,
throughput, and goodput.  Reliability metrics count invalid actions, recovery,
unnecessary escalation, provenance violations, unverified recommendations,
planner latency, and LLM tokens.  Primary curves use cumulative GPU-hours,
with paired nonparametric tests across matched scenarios and all failed
episodes retained.
